
\subsubsection{4.1 Datasets and Models}

\textbf{CIFAR-10 Attribution Benchmark.} CIFAR-10 with 5,000 training samples (subset for tractable LOO computation) and 100 test samples for attribution evaluation.

\textbf{ResNet-18 (Non-Convex Setting).} Modified for CIFAR-10 (conv1 kernel=3, no maxpool), trained for 200 epochs with SGD (lr=0.1, momentum=0.9, weight\_decay=5e-4).

\textbf{Logistic Regression (Convex Setting).} Features: 512-dimensional penultimate layer activations from trained ResNet-18. Regularization: L2 with C=100, solver=lbfgs.

\subsubsection{4.2 Attribution Methods}

Four representative methods spanning distinct computational paradigms were compared: TRAK (random projection), TracIn (gradient similarity across checkpoints), IF (Hessian-weighted gradient similarity), and FastIF (last-layer Arnoldi approximation). All methods were evaluated at matched compute budgets [10, 25, 50, 75, 100] gradient-equivalent operations.

\textbf{Rationale for Method Selection.} The experiments focus on foundational methods representing major paradigm families in data attribution: random projection (TRAK), checkpoint-based gradient similarity (TracIn), full-model Hessian approximation (IF), and efficient last-layer methods (FastIF). This selection enables paradigm comparisons. More recent methods like DataInf and MAGIC target specific architectures (LoRA fine-tuning, metadifferentiation) and would require additional architectural assumptions. Extension to architecture-specific methods is a direction for future work.

\subsubsection{4.3 Implementation Details}

The implementations use gradient-based proxies for computational tractability:

\begin{itemize}
\item \textbf{Ground truth:} FC-layer gradient similarity serves as a proxy for full LOO retraining. This approximation is used in the literature and is sufficient for demonstrating the existence of trade-offs, though it may not match true LOO values.
\item \textbf{Method implementations:} Core algorithmic components are implemented rather than using official library versions to ensure consistent compute normalization across methods.
\end{itemize}

While official implementations (e.g., the TRAK library, Captum for IF) may yield different absolute metric values, the key finding---that structural trade-offs exist between quality dimensions---is expected to be robust to implementation details, as the trade-offs arise from the non-convex geometry of the loss landscape. This limitation is discussed further in Section 6.

\subsubsection{4.4 Statistical Analysis}

\begin{itemize}
\item Bootstrap confidence intervals: 1,000 resamples, 95\% confidence level
\item Method seeds: 3 per configuration
\item LOO seeds: 10 for ground truth variance estimation
\item Total configurations per hypothesis: 60 (4 methods x 5 budgets x 3 seeds)
\end{itemize}
